\documentclass[11pt]{article}
\usepackage[margin=1in]{geometry}
\usepackage[T1]{fontenc}
\usepackage{lmodern,booktabs,amsmath,graphicx}
\usepackage[hidelinks]{hyperref}
\usepackage{xurl}
\setlength{\parindent}{0pt}
\setlength{\parskip}{6pt}
\setlength{\emergencystretch}{3em}
\title{Symmetric Face Support: Tracking Gains and Startup Limits}
\author{Pan Dynamics Collision Avoidance Research}
\date{October 3, 2026}
\begin{document}
\maketitle
\section{Results and comparison protocol}
The selected change is an optional LiDAR rectangle placement rule. Actual
YOLO26x detection, native GNSS/IMU preparation and the unchanged NRMM runtime
were replayed on the previous parked-free capture and a prospectively specified
new capture. Both compare the same input streams, declared target dimensions,
confidence association, causal heading estimation and observer gains. Only
rectangle placement differs between each ordinary edge/symmetric pair.

\begin{center}\small
\begin{tabular}{@{}llrrrr@{}}\toprule
Capture & Placement & Raw pos. & Est. pos. & Est. vel. & Valid \\
 & & (m) & (m) & (m/s) & /300 \\\midrule
Development, 15 m & Edge & 0.1629 & 0.1831 & 0.4195 & 300 \\
 & Symmetric & 0.0674 & 0.0946 & 0.3049 & 300 \\
 & Symmetric + blackout & 0.0657 & 0.1034 & 0.3221 & 285 \\
Confirmation, 20 m & Edge & 0.1513 & 0.1983 & 0.6891 & 300 \\
 & Symmetric & 0.0661 & 0.1430 & 0.4997 & 300 \\
 & Symmetric + blackout & 0.0674 & 0.1725 & 0.5882 & 285 \\
\bottomrule\end{tabular}\end{center}

These are planar vector RMSE over 2.00--5.98 s, matching the previous
experiment's fixed acquisition exclusion. Raw RMSE includes available
measurements; estimator RMSE includes every frame, including dropouts.
All-frame RMSE, p95, maxima, availability and prediction hashes are also
exported in the compact metrics. The development baseline cached replay
reproduced the previous actual detector positions within $1.8\times10^{-14}$ m;
actual inference reproduced the selected prototype within
$1.6\times10^{-14}$ m. Confirmation cases all reran the actual detector.
The previous data guided selection; it is not independent validation.

The improvement is conditional on the predeclared 2 s acquisition exclusion.
Over all 300 confirmation frames, symmetric mode is worse: raw position
RMSE changes from 0.5349 to 0.5548 m, estimated position from 0.6437 to
0.6948 m, and velocity from 2.4888 to 2.6381 m/s. The first frame has
2.167/2.261 m position error and an approximately right-angle fitted-heading
ambiguity. Errors in initial geometry and the causal motion heading interact
before acquisition settles. The change does not solve cold start and must not
be described as a uniform improvement. Over all development frames, estimated
position/velocity do improve from 0.2131 m / 0.9137 m/s to
0.1314 m / 0.8195 m/s. Both statistical views are retained.


\section{Mechanism and implementation}
The old development error after 2 s has mean longitudinal/lateral components
of 0.0527/0.1266 m in CARLA ego axes. Lateral standard deviation is 0.0874 m.
The fitted upper rear surface is narrower than the declared 1.9 m bounding
box. The previous placement anchors one observed endpoint to a box edge,
introducing a lateral offset even when the visible surface is centered.
A point centroid would instead depend strongly on return density.

For each candidate heading, let $\ell_j,u_j$ be the observed support limits
along dimension $d_j$, and $r_j=(u_j-\ell_j)/d_j$. In symmetric mode, select
the axis with largest $r_j$. If $r_j\ge0.8$, place its center at
$(\ell_j+u_j)/2$ and retain face anchoring on the other axis. Otherwise retain
both legacy edge candidates. Initial support limits use the existing 2nd/98th
percentiles; the existing inlier refinement uses extrema. Both fitting stages
apply the same rule. The containment objective and heading prior are unchanged.

This assumes approximately symmetric visibility along the selected axis.
The 0.8 gate tests observed coverage, not the truth of that assumption.
Asymmetric occlusion can still bias the midpoint, so the new behavior is
explicitly enabled with \texttt{--rectangle-placement symmetric}; legacy
\texttt{edge} remains the default. The support fraction is exposed as
\texttt{--rectangle-symmetry-minimum-support}, default 0.8. Outputs record
mode, threshold and selected symmetry axis. No ground-truth bias subtraction,
learned offset, future-frame smoothing or target identity information is used.

The development mean error becomes 0.0543/0.0231 m. Fitted-heading RMSE
falls from 1.107 to 0.680 degrees. Residual longitudinal offset is retained;
approximate vehicle geometry and which body surfaces return LiDAR points
remain relevant. This experiment does not calibrate away that residual.

\section{Fresh confirmation and missing observations}
Before new capture, source hashes, parameters and scoring rules were frozen in
\path{confirmation_plan.json}. The initial gap changed from 15 to 20 m;
LiDAR/IMU seed changed from 20261002 to 20261003, and GNSS seed from
20261003 to 20261004. The target body, map-selected Town10HD\_Opt straight,
spawn index 105, nominal 11 m/s follower speeds and sensor layout stayed the
same. This is a new range/noise realization on the same route, not broad
scenario validation. No prediction parameters were retuned after capture.

An owned CARLA 0.10 instance on port 2011 used the existing scene configuration
to hide 143 static vehicle mesh components before spawning the two actors.
It produced 50 stationary calibration, 150 driving warm-up and 300 evaluation
groups at 0.02 s. All 500 native RGB/LiDAR/GNSS/IMU frame IDs matched and
simulation frames were consecutive. Capture and cleanup details are in
\path{server_lifecycle.json}; shared port 2000 was not contacted or modified.

The extra case replaces only derived RGB paths for frames 150--164 with a
black image, a 0.30 s interruption at 3.00 s. Both captures have exactly 15
missing measurements and reacquire at frame 165. All 300 estimates remain
finite. Normal and interrupted outputs are identical through frame 149,
and all cases have identical ego estimates. Prediction hashes were written
before each separate scorer opened the truth stream. Production perception
summaries contain zero truth comparisons. Calibration and preparation use
native sensor callbacks, fixed georeference and extrinsics only.

\section{Alternatives examined and remaining velocity error}
Using the first 0.4 s of measured target world positions to reset the target
velocity causally at frame 20 reduced development velocity RMSE only from
0.3049 to 0.2934 m/s with a linear fit, or 0.2919 m/s with a quadratic fit.
Position RMSE was 0.0933 and 0.0928 m. The quadratic fit inferred noisy lateral
velocity/acceleration. This limited improvement does not justify introducing
an additional reset into the deployed path; the experiment remains a local
prototype, with no production estimator or adapter change.

A separately labeled oracle diagnostic supplied exact truth-derived target
relative positions while retaining native ego signals and the same estimator.
Its position/velocity RMSE was 0.0389 m / 0.2667 m/s. This is diagnostic only,
not achievable sensor-pipeline performance. It shows that correcting target
geometry alone cannot remove all remaining error in this trajectory.
With symmetric measured geometry, velocity RMSE is 0.5803 m/s in 2--3 s,
0.1403 in 3--4 s, and 0.0879 in 4--6 s. Dynamic response and native ego signal
quality deserve separate experiments; they were not retuned here. The ego
estimated velocity RMSE is 0.0477 m/s (derived native GNSS velocity: 0.1127
m/s). The target accelerates from 4.76 m/s at acquisition to 10.05 m/s at
2 s and 11.12 m/s at 4 s. This prioritizes target transient response as a
follow-up investigation; it does not uniquely identify every error source.

\section{Validation, limits and use}
All 23 Python association, causality and rectangle behavior tests pass. New
checks cover a narrow rear body, side-face localization, rigid transforms,
uneven point density, weak-support fallback, full rectangles, invalid support
fractions, the legacy default and prefix invariance through the actual fitter.
Production estimator/configuration hashes remain unchanged; MATLAB replays
exercise the existing runtime on all 300 frames per case. Complete commands,
versions, timings, audits and diagnostic provenance accompany this report.

Successful RMSE checks do not certify a deterministic error bound. Counts of
raw errors exceeding the configured 0.12 m bound, estimated-domain violations
and actual speed ranges are reported. Some speeds are below the configured
10 m/s target-domain minimum. Native compass heading retains the preceding
experiment's ideal-heading limitation. Runtime bound availability does not
validate these premises. Offline YOLO/LiDAR processing exceeds the six-second
record duration; no 50 Hz real-time or closed-loop safety claim is made.

To enable the selected configuration, retain the previous native pipeline and
use the following options (normal confidence association suffices in the
parked-free scene):
\begin{verbatim}
--camera-id front --causal --target-association confidence
--center-mode rectangle --rectangle-heading-step 1
--heading-iterations 3 --min-depth 0.1 --min-object-height 0
--rectangle-placement symmetric
--rectangle-symmetry-minimum-support 0.8
\end{verbatim}
Raw captures, drivers and generated figures remain ignored local artifacts;
the external archive retains identified exports and source hashes. Tracked
artifacts are the perception change, behavior tests, this English report and
compact results. Unrelated workspace work is excluded.
\IfFileExists{accuracy_trace.pdf}{\begin{figure}[p]\centering
\includegraphics[width=\textwidth]{accuracy_trace.pdf}
\caption{Frozen error traces including acquisition. The vertical line marks
2 s; the green band applies only to the camera-blackout case.}
\end{figure}}{}
\end{document}
